% 2026年9月月度报告；XeLaTeX 编译两遍。
\documentclass[stage2,fontset=fandol]{ydzjreport}
\stagetwosetup
  {阶段二 · 研究推进与成果比较}
  {冯旭东}
  {流道形态}
  {2026-09-01 至 2026-09-30}
  {四种流道；网格无关性；水物性UDF；分区冷却}
  {整理四种流道39个有效工况，完成网格无关性检验并形成约150万网格的后续基准，识别原水物性误差并完成新UDF编译。}
\begin{document}
\makestagetwohead

\section{本月研究任务}
\subsection{任务目标}
作为流道形态设计的主要承担成员，本月将前期多种结构探索收敛到四种代表性形态，服务于申报书提出的分区复合微通道设计\cite{proposal}。
\begin{enumerate}
\item 整理直线、正弦、NACA2416和gao2.1的模型及结果，明确热流密度、总流量与结果保存。
\item 完成四种流道的网格无关性检验，确定后续网格划分的参考规模；修正水物性函数并编译新的UDF。
\item 以最高温度、平均温度、压差和水力功率共同评价性能，给出有数值依据的优劣比较。
\item 将单元特点对应到低温区、变化区和高温区的需求，形成后续复合冷板的候选配置。
\end{enumerate}

\subsection{完成情况}
本月已有较完整的单元仿真资料，整体分区冷板尚未完成，任务状态见表~\ref{tab:tasks}。
\begin{table}[H]
\centering
\caption{本月任务完成情况}
\label{tab:tasks}
\begin{tabularx}{\textwidth}{@{}Y c Y@{}}
\toprule
\rowcolor{ydzjlight} 任务 & 状态 & 说明 \\
\midrule
四种流道模型与性能数据整理 & 已完成 & 形成39个有效工况及比较图表 \\
四种流道网格无关性检验 & 已完成 & 约150万网格及以上表现较理想，作为后续划分基准 \\
水物性UDF修正与编译 & 已完成 & 已识别原函数误差，完成新UDF编译 \\
建立单元性能排序和分区候选配置 & 已完成 & 已比较温度与水力代价，保留条件限制 \\
统一物性及边界条件后的结构定型 & 部分完成 & 新UDF已编译，使用新物性的统一工况复算仍需推进 \\
整板组合与实物测试 & 未开始 & 需先确认单元性能及供液方案 \\
\bottomrule
\end{tabularx}
\end{table}

\section{进阶阅读与理论支撑}
\subsection{本阶段文献清单}
性能分析主要依据已整理的工作簿\cite{data}与分析报告\cite{analysis}，结合物性参考、收敛监测文档及新编译的水物性UDF源码\cite{udf}完善计算基础。资料与对应问题见表~\ref{tab:lit}。
\begin{table}[H]
\centering
\caption{本阶段文献与资料}
\label{tab:lit}
\begin{tabularx}{\textwidth}{@{}c Y Y@{}}
\toprule
\rowcolor{ydzjlight} 序号 & 文献 & 解决了我哪个问题 \\
\midrule
1 & 四种流道数据工作簿与分析报告\cite{data,analysis} & 确定有效工况、指标口径和已知比较限制 \\
2 & IAPWS液态水物性补充发布\cite{iapws} & 给出物性公式的独立校验依据 \\
3 & Fluent收敛监测技术文档\cite{monitor} & 明确温度、压差历史与残差应共同检查 \\
4 & 新水物性UDF源码\cite{udf} & 明确温度相关物性、适用温区及比热与焓的一致性 \\
\bottomrule
\end{tabularx}
\end{table}

\subsection{对本方向的启示}
物性参考和收敛监测要求，使评价重点由单点低温结果转向可比较的性能证据。数据分析不仅需要看热流密度，还要核对总流量、加热面积与实际时刻；求解稳定性不仅需要看残差，还应观察温度与压差随时间的变化。本月已推进物性修正并编译新UDF，后续将在此基础上统一比较条件，避免将不同模型配置的差异全部归因于形态。

最高温度和体积平均温度分别反映局部热点与总体热状态。水力功率按$P_{\mathrm h}=\Delta p\,Q_{\mathrm v}$估算，峰值热阻按
\[
 R_{\max}=\frac{T_{\max}-T_{\mathrm{in}}}{qA_{\mathrm h}}
\]
计算。归一化有助于比较，但不能消除几何、流量和瞬态储热差异。

\section{研究进展}
\subsection{已完成的工作}
\begin{enumerate}
\item 整理直线通道、正弦曲线通道、NACA2416翼型特征流道与gao2.1自定义流道。gao2.1继续按现有模型编号归档，不将其与前期其他肋流道混称。
\item 汇总39个有效工况的冷板固体最高温度、体积平均温度、进出口压差与水力功率，并建立随热流密度变化的比较资料。
\item 核对共有工况，分别比较一般负荷区的三方案与500 W/cm$^2$时的两方案，不外推不存在的工况。
\item 对直线、sin曲线、NACA2416和gao2.1分别开展网格无关性检验；对不同网格密度的计算结果进行比较，发现网格数量达到约150万及以上时表现较为理想。
\item 仿真后核查前人使用的水物性UDF，发现其存在误差；完成新水物性UDF的编译，为后续提高计算可信度建立新的物性基础。
\end{enumerate}

\textbf{网格无关性检验与后续基准。} 四种流道均已完成网格无关性检验，本阶段将约$1.5\times10^6$个网格单元作为后续划分的参考规模。以后建模时以这一规模为起点，并在窄隙、尖角和流固界面局部加密；几何或工况发生变化时仍需复核网格敏感性。该规模是本轮检验形成的基准，不表示达到该数量即可自动满足所有模型的精度要求。

\textbf{水物性UDF修正进展。} 本月已编译新的水物性UDF\cite{udf}。源码注明采用101325 Pa下REFPROP WATER.FLD v9.1液态水数据拟合，提供随温度变化的密度、黏度、导热系数和比热；比热与显焓采用同一多项式及其积分，保证二者定义一致。源码对流体材料温度280--360 K进行范围检查，并提供300 K和330 K的物性自检输出。此次进展由发现原函数误差推进到新函数编译，后续还需完成挂接、物性校验与代表工况复算。

四方案既有39个工况均为铜冷板、水工质、入口速度2 m/s、入口温度300 K（26.85 $^\circ$C），采用单相、层流、瞬态模型；实际参数见表~\ref{tab:conditions}。下文性能数值沿用这些已有结果，未将其标作新UDF复算结果。申报书中的湍流模型属于拟研究方案，不能替代已保存模型的实际配置。
\begin{table}[H]
\centering
\caption{四种流道的工况范围与比较条件}
\label{tab:conditions}
\begin{tabularx}{\textwidth}{@{}Y c Y Y Y Y@{}}
\toprule
\rowcolor{ydzjlight} 流道 & 点数 & 热流密度\newline W/cm$^2$ & 总流量\newline L/min & 加热面积\newline cm$^2$ & 时刻\newline s \\
\midrule
直线 & 10 & 100--325 & 0.96 & 0.1050 & 0.60 \\
正弦 & 10 & 100--325 & 0.72 & 0.1225 & 0.80 \\
NACA2416 & 7 & 150--500 & 0.96 & 0.1225 & 0.80 \\
gao2.1 & 12 & 500--1050 & 0.54 & 0.1225 & 0.80 \\
\bottomrule
\end{tabularx}
\end{table}


\subsection{关键结果}
代表性结果见表~\ref{tab:results}。水力功率仅估算冷板静压损失的代价，不含泵效率和系统其他损失。
\begin{table}[H]
\centering
\caption{既有工况在共同热流密度下的温度与水力性能}
\label{tab:results}
\begin{tabularx}{\textwidth}{@{}c Y Y Y Y Y@{}}
\toprule
\rowcolor{ydzjlight} $q$ / W/cm$^2$ & 流道 & 最高温度 / $^\circ$C & 平均温度 / $^\circ$C & 压差 / kPa & 水力功率 / W \\
\midrule
300 & 直线 & 52.50 & 38.43 & 3.074 & 0.04919 \\
300 & 正弦 & 62.27 & 44.84 & 4.914 & 0.05897 \\
300 & NACA2416 & 72.48 & 45.99 & 0.979 & 0.01566 \\
500 & gao2.1 & 68.50 & 42.55 & 12.457 & 0.11211 \\
500 & NACA2416 & 99.60 & 56.82 & 0.926 & 0.01482 \\
\bottomrule
\end{tabularx}
\end{table}

\textbf{一般负荷下的比较。} 在150、200、250、300 W/cm$^2$四个共同工况，最高温度与平均温度均为直线低于正弦、正弦低于NACA2416；压差与水力功率均为NACA2416低于直线、直线低于正弦。

在300 W/cm$^2$时，直线最高温度比正弦低9.77 $^\circ$C，比NACA2416低19.98 $^\circ$C；相对正弦，压差低约37.4\%，水力功率低约16.6\%。直线因此是当前温度控制优先区域的基础候选。NACA2416的压差和水力功率均比直线低约68.2\%，但最高温度较高，适合有充分温度余量且节省泵耗优先的区域。正弦当前未表现出优先布置优势。

这里的温度差是现有方案在相同热流密度标签下的结果差异，不表示同一冷板改造前后的实测温降。300 W/cm$^2$时，直线总加热功率为31.50 W，正弦与NACA2416为36.75 W；直线加热面积小约14.3\%，且正弦流量较小，因此差异仍包含工况影响。

\textbf{高负荷下的比较。} 在500 W/cm$^2$时，gao2.1最高温度比NACA2416低31.10 $^\circ$C，平均温度低14.27 $^\circ$C；峰值热阻分别约为0.680 K/W和1.188 K/W，gao2.1低约42.8\%。同时，其压差约为NACA2416的13.45倍，水力功率约为7.57倍，说明热点控制改善伴随更高的流动代价。两个结果的实际时刻分别为0.80 s与0.41 s，需统一后复核。

gao2.1在500--1050 W/cm$^2$范围内，最高温度由68.50 $^\circ$C升至110.36 $^\circ$C。1050 W/cm$^2$时平均温度为57.50 $^\circ$C，热点高出平均温度52.86 K，仅看平均值会掩盖局部高温。该方案是高负荷区的优先复算候选，但缺少与直线、正弦的高负荷共同数据，不能认定四方案中的全面最优。

\subsection{与预期的差异}
结构复杂并不自动带来更好的综合表现。正弦在当前配置下同时表现出较高温度与较高水力代价，说明曲线形态的潜在强化作用尚未转化为可确认的优势；NACA2416的突出特点是低流阻，而不是最低热点温度。gao2.1在高负荷下有较低温度，却有更大的压降。

这些差异提示，分区设计不应把“复杂流道”简单对应到“散热更强”。后续应通过控制变量试算，将形态作用与受热面积、流量、物性和计算时刻的影响区分开，再判断是否存在可用于整板的收益。

\section{问题与对策}
\subsection{遇到的主要问题}
\begin{enumerate}
\item 水物性存在一定偏差。在298.15 K校验点，密度、比热和黏度相对参考值分别偏低约60.8\%、10.5\%和26.7\%，可能影响绝对性能及排序。

\item 单块性能不能直接代表整板性能。并联时高阻力支路可能供液不足，串联时上游吸热会抬高下游入口水温。
\end{enumerate}

\subsection{已尝试的方案}
\begin{enumerate}
\item 已核查原水物性函数误差，并编译新的水物性UDF；新源码包含温度范围检查、物性自检及比热--焓一致处理。后续需以新函数重新计算，确认修正对温度与压差的影响。
\item 已完成四种流道的网格无关性检验，形成约150万网格的后续划分基准；各结构的网格规模、关键物理量及相邻网格结果变化需同步归档。
\item 已统一单位和指标定义，并限制比较到共同热流密度，同时计算峰值热阻与水力功率；这些处理仍不能消除几何与时刻差异。
\item 已将单元选型整理为条件性配置建议，保留各区温度余量和泵功约束；尚未用独立单块排序替代整板分流验证。
\end{enumerate}

\subsection{下一步对策}
\begin{enumerate}
\item 挂接已编译的新UDF，核查入口及实际流体温区，完成物性自检并重新求解；超过280--360 K时先扩展拟合范围，不能按比例缩放现有温度或压差。
\item 统一受热面积、总流量和比较时刻，补齐四方案300及500 W/cm$^2$对照，检查温度、压差和热量收支历史。
\item 构建分区复合冷板，验证支路流量与区域入口水温，补充温度标准差和区域间温差。
\item 设计并制作各小块连接处的过渡与分流通道，通过优化连接段截面及分流布局，减少局部回流、流动突变和偏流，使各小块入口的速度分布尽量均匀、流动整体保持稳定；结合入口速度分布、支路流量及连接段压降评价效果，并按各区冷却需求合理分配流量。
\end{enumerate}

\section{阶段性结论}
\subsection{可以确定的}
\begin{enumerate}
\item 已形成四种形态、39个有效工况的数据和图表，单元层面的结果资料较为完整，结构工作已由多方案探索推进到重点形态筛选。
\item 已完成四种流道的网格无关性检验，形成约150万网格的后续基准；已识别原水物性UDF误差并完成新UDF编译，计算基础进一步完善。
\item 现有共同工况显示，直线温度较低，NACA2416水力代价较低，正弦暂缺优先配置依据；gao2.1提供了高负荷热点控制的验证方向。
\item 单元形态应按局部需求配置，而不是全板采用同一种结构，候选方案见表~\ref{tab:zones}。
\end{enumerate}

\begin{table}[H]
\centering
\caption{流道形态与温区需求的候选对应}
\label{tab:zones}
\begin{tabularx}{\textwidth}{@{}c Y Y@{}}
\toprule
\rowcolor{ydzjlight} 分区 & 优先候选 & 选用条件与证据 \\
\midrule
低温区 & NACA2416，或直线 & 有温度余量时选低阻力，否则选直线 \\
变化区 & NACA2416为基础，结合供液调节 & 共同工况控温较好；动态适应性待验证 \\
高温区 & gao2.1 & 高负荷候选；需满足温度和泵功约束 \\
后续对照 & 正弦 & 现有工况不优先；统一条件后再评价 \\
\bottomrule
\end{tabularx}
\end{table}

\subsection{尚不能确定的}
\begin{enumerate}
\item 修正物性并统一条件后，现有排序是否保持？需要重新计算而非仅重算后处理指标。
\item 各小块的温度边界与最终布局尚不能确定。三个温区为相对热需求分区，应结合允许最高温度、温度余量、热流及温度梯度确定，不能将热流密度当作温度。
\item 整板均温和节能效果是否优于单一流道？需要组合模型、时间序列和后续实测证据；目前尚不能声称实物已验证。
\end{enumerate}

\section{下月计划}
2026年10月将优先修正比较基础，再推进组合验证，计划见表~\ref{tab:plan}。
\begin{table}[H]
\centering
\caption{下月工作计划}
\label{tab:plan}
\begin{tabularx}{\textwidth}{@{}Y Y c@{}}
\toprule
\rowcolor{ydzjlight} 事项 & 预期产出 & 时间 \\
\midrule
挂接并校验新UDF，统一总流量、受热面积与比较时刻 & 物性自检记录与复算工况清单 & 第1周 \\
设计小块连接处的过渡与分流通道，改善入口流动稳定性与均匀性 & 连接段形态方案及入口流场对比 & 第2--3周 \\
构建候选复合冷板，核对分流与区域间温差 & 分区布局及动态调控基础数据 & 第4周 \\

\bottomrule
\end{tabularx}
\end{table}

\begin{thebibliography}{9}
\bibitem{proposal} 王佳蔚等. 大创项目申报书初稿，研究内容与实施方案. 西安交通大学，2026.
\bibitem{data} 四种流道\_数据与折线图.xlsx. 结论、共同工况、工况核对及补充与来源工作表.
\bibitem{analysis} 四种流道\_分析报告.pdf. 第1--7页，阶段数据整理与复算建议. 2026-10-07.
\bibitem{iapws} IAPWS. SR6-08(2011)：0.1 MPa下液态水物性补充发布. 2011. \href{https://iapws.org/technical-guidance/release/LiquidWater}{官方发布页面}.
\bibitem{monitor} Ansys. Fluent用户指南：收敛监测. 2025 R1. \href{https://ansyshelp.ansys.com/public/Views/Secured/corp/v251/en/flu_ug/flu_ug_sec_solve_monitor.html}{官方技术文档}.
\bibitem{udf} 项目组. water\_properties\_refprop\_fit .c：液态水物性拟合与Fluent UDF源码. 项目计算资料，2026.
\end{thebibliography}
\end{document}
